\subsubsection{A KL Projection View of our Reward}
In this subsection we analyze PivotRL from a KL projection view that isolates the effect of the matching-style reward
$r_{\name}(s,a)$ in the RL objective.
The analysis only depends on the induced matching set, so it applies to any matching rule that can be written in this indicator form.
The key message is that the population-loss optimizer minimally perturbs the reference policy among all policies that achieve the same expert-matching level.
Because this reward depends only on whether an action falls inside the matching set, the optimizer changes the reference only through a set-level reallocation of mass, while preserving the relative probabilities within the matching set and within its complement.

\begin{theorem}[RL with $r_{\name}$ as KL projection under matching-style rewards]\label{thm:pivotrl_kl_projection}
Assume the action space $\mathcal A$ is finite and $\pi_0(a\mid s)>0$ for all $s$ and $a$.
Fix $\beta>0$ and $\epsilon\geq 0$.
Consider the RL loss
\begin{align*}
\mathcal L_{\text{PivotRL},\beta}(\pi)
=
\mathbb E_{s\sim\mathcal D_{\sft}}\Big[
-\mathbb E_{a\sim\pi(\cdot\mid s)}[r_{\name}(s,a)]
+ \beta \cdot \KL(\pi(\cdot\mid s)\|\pi_0(\cdot\mid s))
\Big]
\end{align*}
with $r_{\name}$ defined in Equation~\eqref{eq:pivotrl-reward}.
For each $s$, let
\begin{align*}
\rho_\epsilon(s)=\pi_0(\mathcal M_\epsilon(s)\mid s),
\qquad
q_{\beta,\epsilon}(s)=\frac{\rho_\epsilon(s)e^{1/\beta}}{(1-\rho_\epsilon(s))+\rho_\epsilon(s)e^{1/\beta}}.
\end{align*}
Then the unique minimizer $\pi_{\beta,\epsilon}^\star$ of $\mathcal L_{\text{PivotRL},\beta}$ (up to $\mathcal D_{\sft}$-a.e. equality) is characterized as follows:
for $\mathcal D_{\sft}$-a.e. $s$, among all distributions $\pi(\cdot\mid s)$ satisfying the matching constraint
\begin{align*}
\pi(\mathcal M_\epsilon(s)\mid s)=q_{\beta,\epsilon}(s),
\end{align*}
$\pi_{\beta,\epsilon}^\star(\cdot\mid s)$ is the unique minimizer of
\begin{align*}
\KL(\pi(\cdot\mid s)\|\pi_0(\cdot\mid s)).
\end{align*}
Equivalently, if $0<\rho_\epsilon(s)<1$, then
\begin{align}\label{eq:teacher-policy}
\pi_{\beta,\epsilon}^\star(a\mid s)
=
\begin{cases}
\frac{q_{\beta,\epsilon}(s)}{\rho_\epsilon(s)}\pi_0(a\mid s), & a\in\mathcal M_\epsilon(s),\\[6pt]
\frac{1-q_{\beta,\epsilon}(s)}{1-\rho_\epsilon(s)}\pi_0(a\mid s), & a\notin\mathcal M_\epsilon(s),
\end{cases}
\end{align}
while if $\rho_\epsilon(s)=1$, then $\pi_{\beta,\epsilon}^\star(\cdot\mid s)=\pi_0(\cdot\mid s)$.
\end{theorem}

Theorem~\ref{thm:pivotrl_kl_projection} characterizes the ``conservatism'' of $r_{\name}$:
for each $s$, the optimal policy $\pi_{\beta,\epsilon}^\star(\cdot\mid s)$ is the unique distribution that achieves the matching level $q_{\beta,\epsilon}(s)$ on the acceptable set $\mathcal M_\epsilon(s)$ and perturbs the reference $\pi_0(\cdot\mid s)$ as little as possible in KL, which is precisely an information projection~\citep{csiszar1975divergence,bregman1967relaxation}.
From Equation~\eqref{eq:teacher-policy}, we also infer that $\pi_{\beta,\epsilon}^\star$ differs from $\pi_0$ only through a two-block reallocation of mass:
\begin{align*}
\frac{\pi_{\beta,\epsilon}^\star(a\mid s)}{\pi_{\beta,\epsilon}^\star(b\mid s)}=\frac{\pi_0(a\mid s)}{\pi_0(b\mid s)}, ~\forall~a,b\in \mathcal M_\epsilon(s),
\end{align*}
and
\begin{align*}
\frac{\pi_{\beta,\epsilon}^\star(a\mid s)}{\pi_{\beta,\epsilon}^\star(b\mid s)}=\frac{\pi_0(a\mid s)}{\pi_0(b\mid s)}, ~\forall~a,b\notin \mathcal M_\epsilon(s).
\end{align*}
This implies that $\pi_{\beta,\epsilon}^\star(\cdot\mid s)$ preserves the reference ranking among all matching actions and, separately, among all non-matching actions.
Thus a matching-style reward increases only the total mass assigned to the acceptable set, rather than forcing all additional mass onto one canonical response.
Together, these perspectives explain why $r_{\name}$ can increase expert matching while maintaining the reference distribution over non-golden actions---and, more generally, within each reward-equivalence class---thereby better preserving out-of-domain performance and mitigating catastrophic forgetting.
The proof is deferred to Appendix~\ref{sec:proof-kl-proj}.